\section*{Hoofdstuk \ref{ch:b1:earthquake-source} --- De aardbevingsbron}

\begin{solution}{exo:b1:earthquake-source:1}
(a) $\log_{10}(\num{4e17}) = 17.60$, dus $M_w = \tfrac23(17.60 - 9.1) = 5.67$.
(b) $\log_{10}(\num{5.59e22}) = 22.75$ en $M_w = \tfrac23(22.75 - 9.1) =
9.10$. (c) $M_0 = 10^{1.5\times7.0 + 9.1} = 10^{19.6} = \qty{3.98e19}{N.m}$.
\end{solution}

\begin{solution}{exo:b1:earthquake-source:2}
(a) $T = D/\dot s = \qty{3}{m}/\qty{0.025}{m.yr^{-1}} = \qty{120}{yr}$.
(b) Met $c = 0.5$, $T = D/(c\dot s) = \qty{240}{yr}$: het kruipen neemt de helft van
de beweging op, en de aardbevingen, met dezelfde verschuiving, komen half zo vaak voor.
\end{solution}

\begin{solution}{exo:b1:earthquake-source:3}
Volgens \cref{prop:b1:earthquake-source:moment-steps} is de verhouding van de momenten
$10^{1.5\times(7.4 - 6.2)} = 10^{1.8} = 63.1$: er zijn 63 aardbevingen van
magnitude 6.2 nodig om het moment van één van magnitude 7.4 vrij te geven.
\end{solution}

\begin{solution}{exo:b1:earthquake-source:4}
Met $b = 1$ deelt elke eenheid magnitude het aantal door 10: 10 aardbevingen van
magnitude 5 of meer per jaar, 1 van magnitude 6 of meer, en 0.1 van magnitude 7 of
meer, dat is gemiddeld één elke 10 jaar.
\end{solution}

\begin{solution}{exo:b1:earthquake-source:5}
(a) $A = \num{4.3e5} \times \qty{1.5e4}{m} = \qty{6.45e9}{m^2}$, en
$M_0 = \num{3e10}\times\num{6.45e9}\times4 = \qty{7.74e20}{N.m}$;
$\log_{10}M_0 = 20.89$ en $M_w = \tfrac23(20.89 - 9.1) = 7.86$. (b) Een diepte van
\qty{10}{km} deelt het oppervlak en het moment door 1.5, wat de magnitude verlaagt met
$\tfrac23\log_{10}1.5 = 0.12$, tot 7.74. (c) De eerste, 7.86, rondt af tot de 7.9 van
de catalogus; de tweede is een beetje laag. De magnitude is niet erg gevoelig voor de
diepte: een factor 1.5 op het oppervlak is slechts 0.12 magnitude.
\end{solution}

\begin{solution}{exo:b1:earthquake-source:6}
(a) $D = M_0/(\mu A) = \num{5.59e22}/(\num{3e10}\times\num{4e5}\times\num{1.5e5})
= \qty{31.1}{m}$. (b) $\qty{31.1}{m}/\qty{0.083}{m.yr^{-1}} = \qty{374}{yr}$.
(c) Met de vastgeklemde fractie 0.8, $31.1/(0.8\times0.083) = \qty{468}{yr}$; met
0.5, \qty{749}{yr}. Hoe minder het grensvlak is vastgeklemd, hoe langer het duurt
voordat de verschuiving van één grote aardbeving is opgebouwd; de gemiddelde verschuiving
die hier met een rechthoek van \qtyproduct{400 x 150}{km} en $\mu = \qty{30}{GPa}$ is
gevonden, is zelf slechts een schatting, en de werkelijke verschuiving was allesbehalve
uniform.
\end{solution}

\begin{solution}{exo:b1:earthquake-source:7}
(a) De \omterm{def:b1:geological-maps:attitude}{strijk} is $\phi = 329$, de \omterm{def:b1:geological-maps:attitude}{helling-richting} $a = 059$, de helling $\delta =
\ang{8}$. Dan is $\mathbf s = (\sin329, \cos329, 0) = (-0.5150, 0.8572, 0)$,
$\mathbf u = (-\cos8\sin59, -\cos8\cos59, \sin8) = (-0.8488, -0.5100, 0.1392)$,
$\mathbf n = (\sin8\sin59, \sin8\cos59, \cos8) = (0.1193, 0.0717, 0.9903)$,
en met $r = \ang{110}$,
\[
  \mathbf d = \cos110\,\mathbf s + \sin110\,\mathbf u = (-0.6215, -0.7724, 0.1308) .
\]
(b) Het \omterm{def:b1:earthquake-source:focal-mechanism}{hulpvlak} heeft de normaal $\mathbf d$ en de verschuiving $\mathbf n$. Naar boven
genomen helt het onder $\arccos 0.1308 = \ang{82.5}$ naar
$\arctan(-0.6215/-0.7724) = 218.8$, dus is de \omterm{def:b1:geological-maps:attitude}{strijk} $128.8$; de verschuiving
$\mathbf n$, geprojecteerd op haar strijk- en opwaartse-helling-vectoren, geeft de \omterm{def:b1:faults-folds-fabrics:fault-elements}{rake}
\ang{87.2}. Afgerond: $129/83/87$, het tweede vlak van de gepubliceerde oplossing.
(c) $\mathbf P = (\mathbf n - \mathbf d)/\sqrt2$ heeft trend $221$ en plunge
\ang{37}; $\mathbf T = (\mathbf n + \mathbf d)/\sqrt2$ heeft trend $036$ en plunge
\ang{52} (gepubliceerd: $222/38$ en $36/52$). (d) Het ondiepe vlak, dat \ang{8} naar
$059$ helt, langs de \omterm{def:b1:active-margins:slab}{subducerende plaatzijde}: een \omterm{def:b1:faults-folds-fabrics:fault-types}{omgekeerde breuk} met een \omterm{def:b1:faults-folds-fabrics:fault-elements}{rake} van \ang{110},
en, omdat het minder dan \ang{45} helt, een \omterm{def:b1:faults-folds-fabrics:fault-types}{overschuiving}, het plaatgrensvlak van de
subductiezone.
\end{solution}

\begin{solution}{exo:b1:earthquake-source:8}
(a) Het vlak helt naar $135$, dus $\mathbf n = (\sin135, \cos135, 0) =
(0.707, -0.707, 0)$; een dextrale verschuiving heeft $r = \ang{180}$, $\mathbf d = -\mathbf s
= (-0.707, -0.707, 0)$. Naar het noorden, $\mathbf x = (0, 1, 0)$:
$(\mathbf n\cdot\mathbf x)(\mathbf d\cdot\mathbf x) = (-0.707)(-0.707) > 0$, een
duw. Oostwaarts: $(0.707)(-0.707) < 0$, een trek. Zuidwaarts: een duw; westwaarts: een trek.
(b) $\mathbf P = (\mathbf n - \mathbf d)/\sqrt2 = (1, 0, 0)$: de \omterm{def:b1:earthquake-source:focal-mechanism}{P-as} is horizontaal en
heeft trend $090$--$270$; $\mathbf T = (0, -1, 0)$ heeft trend $000$--$180$. (c) Het
verticale vlak dat $135$--$315$ strijkt, heeft de normaal $\mathbf d$. Met de
verschuiving $\mathbf n$ beweegt het blok aan zijn zuidwestelijke kant, waarheen
$\mathbf d$ wijst, naar $135$; gezien vanuit het noordoostelijke blok, naar het
zuidwesten kijkend, is dat naar links: een sinistrale verschuiving. Volgens
\cref{prop:b1:earthquake-source:auxiliary} geeft deze verschuiving dezelfde \omterm{def:b1:earthquake-source:focal-mechanism}{eerste
bewegingen} in elke richting.
\end{solution}

\begin{solution}{exo:b1:earthquake-source:9}
(a) $V = \qty{40}{mm.yr^{-1}}$ en $v = V/4$ op \qty{10}{km}, dus volgens
\cref{prop:b1:earthquake-source:profile}(b) is $D = \qty{10}{km}$; direct,
$D = x/\tan(\pi v/V) = 10/\tan(\pi/4) = \qty{10}{km}$. (b) $v(30) =
(40/\pi)\arctan 3 = \qty{15.9}{mm.yr^{-1}}$. (c) $V/(\pi D) =
\num{0.040}/(\pi\times\num{10000}) = \num{1.27e-6}$ per jaar.
\end{solution}

\begin{solution}{exo:b1:earthquake-source:10}
(a) De momenten van de klassen zijn nog steeds $C q^k$ met $q = 10^{1.5 - b}$, maar nu
geldt $q < 1$: het moment van de klassen neemt af naarmate de magnitudes groter worden,
en de onderste klasse, $k = 0$, draagt het aandeel $C/S = (1 - q)/(1 - q^{K+1}) > 1 - q$,
het spiegelbeeld van \cref{prop:b1:earthquake-source:largest}. De kleinste geteldde
aardbevingen zouden dan domineren, en het totaal zou zonder grens groeien naarmate steeds
kleinere werden meegeteld. (b) $1 - 10^{-0.7} = 0.80$: ten minste \qty{80}{\percent}.
(c) $\log_{10}N = 8.20 - 1.01\times9 = -0.89$, dus $N = 0.13$ per jaar en ongeveer 6.4
in vijftig jaar, drie keer de twee waargenomen. De lijn, gefit tussen magnitudes 6 en
7.5, overschat de grootste aardbevingen; ofwel buigt de wet aan de top naar beneden, zoals
de eindige grootte van breuken suggereert, ofwel zijn vijftig jaar te weinig om
gebeurtenissen te bemonsteren die eens in de decennia komen. Beide zijn mogelijk met twee
gebeurtenissen.
\end{solution}

\begin{solution}{exo:b1:earthquake-source:11}
(a) In 150 jaar verplaatst een paal op $x$ zich met $150\,v(x) = 150\times(34/\pi)
\arctan(x/15)$ millimeter ten opzichte van de paal op de breuk: \qty{0.11}{m} op
\qty{1}{km}, \qty{1.28}{m} op \qty{15}{km}, \qty{2.08}{m} op \qty{50}{km}, naar het zuiden
voor een rechtsdraaiende breuk gezien vanuit het westen. (b) Het verschuivingstekort is
$s = \qty{0.034}{m.yr^{-1}}\times\qty{150}{yr} = \qty{5.1}{m}$. Tijdens de aardbeving
verplaatsen de palen zich over $(s/\pi)\arctan(D/x)$: \qty{2.44}{m} op \qty{1}{km},
\qty{1.28}{m} op \qty{15}{km}, \qty{0.47}{m} op \qty{50}{km}. (c) De sommen zijn
$0.11 + 2.44$, $1.28 + 1.28$ en $2.08 + 0.47$, alle \qty{2.55}{m} $= s/2$: na de
aardbeving is het hek aan de oostkant weer recht, verplaatst met $s/2$, en de westkant,
uit symmetrie, met $-s/2$.
\end{solution}

\begin{solution}{exo:b1:earthquake-source:12}
(a) $\mathbf P = (0, 1, 0)$ en $\mathbf T = (1, 0, 0)$. Het paar
$\mathbf n = (\mathbf T + \mathbf P)/\sqrt2 = (0.707, 0.707, 0)$,
$\mathbf d = (\mathbf T - \mathbf P)/\sqrt2 = (0.707, -0.707, 0)$ voldoet aan
$\mathbf P = (\mathbf n - \mathbf d)/\sqrt2$ en $\mathbf T = (\mathbf n +
\mathbf d)/\sqrt2$. Beide vectoren zijn horizontaal, dus zijn beide \omterm{def:b1:earthquake-source:focal-mechanism}{knoopvlakken}
verticaal: het vlak met normaal $\mathbf n$ (naar $045$) strijkt $135$--$315$, het vlak
met normaal $\mathbf d$ (naar $135$) strijkt $045$--$225$. (b) Op het vlak dat $315$
strijkt, beweegt het blok aan de noordoostkant langs $\mathbf d$, naar $135$: gezien
vanuit het andere blok, naar rechts, een rechtsdraaiende verschuiving. Op het vlak dat
$045$ strijkt, beweegt het blok aan de zuidoostkant langs $\mathbf n$, naar $045$: een
linksdraaiende verschuiving. (c) Een nieuwe breuk maakt $\ang{45} - \phi/2 =
\ang{30}$ met $\sigma_1$: met $\sigma_1$ met trend $000$ zouden de breuken $030$
(\omterm{def:b1:faults-folds-fabrics:fault-types}{sinistraal}) en $330$ (\omterm{def:b1:faults-folds-fabrics:fault-types}{dextraal}) \omterm{def:b1:geological-maps:attitude}{strijken}, \ang{15} van de \omterm{def:b1:earthquake-source:focal-mechanism}{knoopvlakken}. De \omterm{def:b1:earthquake-source:focal-mechanism}{P-as} ligt
door de meetkunde op \ang{45} van de \omterm{def:b1:earthquake-source:focal-mechanism}{knoopvlakken}, wat de wrijving van het gesteente ook
is, dus valt zij alleen samen met $\sigma_1$ als de breuk onder \ang{45} ervan is
geschoven. (d) De breuk die $040$ strijkt, ligt \ang{5} van het \omterm{def:b1:earthquake-source:focal-mechanism}{knoopvlak} dat $045$
strijkt: dat vlak is waarschijnlijk de breuk, en de verschuiving erop is
linksdraaiend.
\end{solution}

\begin{solution}{pb:b1:earthquake-source:1}
\textbf{1.} De intervallen zijn $1290 - 1050 = 240$, $1480 - 1290 = 190$ en
$1712 - 1480 = 232$ jaar; hun gemiddelde is \qty{220.7}{yr}.

\textbf{2.} Het gemiddelde verspringing is $(7.0 + 6.0 + 5.5 + 7.5)/4 = \qty{6.5}{m}$, en
$T = 6.5/0.030 = \qty{216.7}{yr}$, dicht bij het gemiddelde interval.

\textbf{3.} De verspringingen van de laatste drie aardbevingen, $6.0 + 5.5 + 7.5 =
\qty{19.0}{m}$, zijn opgebouwd tussen 1050 en 1712, in 662 jaar: $\qty{19.0}{m}
/\qty{662}{yr} = \qty{28.7}{mm.yr^{-1}}$, \qty{4}{\percent} minder dan de
geodetische \qty{30}{mm.yr^{-1}}. De twee komen overeen binnen de onzekerheid van de
dateringen en verspringingen: het segment geeft bijna zijn hele belasting in zijn grote
aardbevingen vrij.

\textbf{4.} $2025 - 1712 = 313$ jaar, en het tekort is
$\qty{0.030}{m.yr^{-1}}\times\qty{313}{yr} = \qty{9.39}{m}$.

\textbf{5.} Het tekort overtreft elke verspringing in de sleuf, de grootste daarvan
\qty{7.5}{m}, en de tijd sinds 1712 overtreft het langste interval, 240 jaar. In die zin
is het segment achterstallig; het betekent niet dat de aardbeving op handen is, alleen dat
de breuk al meer verschuiving bevat dan zij ooit in één keer heeft vrijgegeven in het
verslag.

\textbf{6.} De intervallen lopen van 190 tot 240 jaar en de verspringingen van
\qtyrange{5.5}{7.5}{m}: de cyclus is niet periodiek, en de sterkte van de breuk bij de
volgende aardbeving is onbekend. De sleuf geeft een gemiddelde en een spreiding, waaruit
een kans kan worden geschat, geen datum.

\textbf{7.} $A = \qty{250}{km}\times\qty{15}{km} = \qty{3750}{km^2} =
\qty{3.75e9}{m^2}$.

\textbf{8.} $M_0 = \num{3e10}\times\num{3.75e9}\times\num{9.39} =
\qty{1.06e21}{N.m}$.

\textbf{9.} $M_0 = \num{3e10}\times\num{3.75e9}\times\num{7.5} =
\qty{8.44e20}{N.m}$, dus $M_w = \tfrac23(20.926 - 9.1) = 7.88$.

\textbf{10.} De verhouding is $9.39/7.5 = 1.25$, wat $\tfrac23\log_{10}1.25 =
0.07$ magnitude is.

\textbf{11.} Half het oppervlak halveert het moment en verlaagt de magnitude met
$\tfrac23\log_{10}2 = 0.20$: $M_w = 7.75$.

\textbf{12.} $\mu A V = \num{3e10}\times\num{3.75e9}\times\num{0.030} =
\qty{3.38e18}{N.m}$ per jaar.

\textbf{13.} $b = 1$. $N(5) = 10^{4.3 - 5} = 0.20$ per jaar en $N(6) = 0.020$ per
jaar.

\textbf{14.} $N(7) = 0.0020$ per jaar: één elke 500 jaar; $N(8) = 0.00020$ per
jaar: één elke \num{5000} jaar.

\textbf{15.} De wet geeft één aardbeving van magnitude 7.95 of meer elke
$1/10^{4.3 - 7.95} = \num{4500}$ jaar, de sleuf één elke 220 jaar, twintig keer vaker.
De wet is gefit op de vele kleine aardbevingen van alle breuken van de regio;
geëxtrapoleerd naar de grote aardbevingen van één segment, onderschat zij ze sterk. Het
\omterm{def:b1:earthquake-source:seismic-cycle}{herhalingsinterval} van grote aardbevingen op een breuk wordt genomen uit haar eigen verslag
en uit haar belasting, niet uit de kleine.

\textbf{16.} $\num{5.5e17}/\num{3.38e18} = 0.16$: de matige aardbevingen geven
\qty{16}{\percent} van het moment vrij dat het segment opbouwt.

\textbf{17.} Volgens \cref{prop:b1:earthquake-source:moment-steps} is het
$10^{1.5\times(7.95 - 6)} = 10^{2.92} \approx 840$ keer meer.

\textbf{18.} Met $b = 1$ wordt het moment van de aardbevingen gedomineerd door de
grootste: wat de matige aardbevingen ook doen, zij geven een klein aandeel vrij, hier
\qty{16}{\percent}, en de rest, \qty{84}{\percent}, het verschuivingstekort, wacht op de
grote aardbeving.

\textbf{19.} $v(-x) = -v(x)$ voor elk paar, dus is het profiel antisymmetrisch. Ver van de
breuk neigen de snelheden naar $\pm V/2 = \pm\qty{15}{mm.yr^{-1}}$.

\textbf{20.} $D = x/\tan(\pi v/V) = 15/\tan(\pi\times7.5/30) = 15/\tan\ang{45} =
\qty{15}{km}$: op de \omterm{def:b1:earthquake-source:locking-depth}{vastklemdiepte} is de snelheid $V/4$.

\textbf{21.} Op \qty{5}{km}: $D = 5/\tan(\pi\times3.1/30) = 5/\tan\ang{18.6} =
\qty{14.9}{km}$; op \qty{40}{km}: $D = 40/\tan\ang{69.6} = \qty{14.9}{km}$. De
snelheden zijn gegeven tot op \qty{0.1}{mm.yr^{-1}}, en die afronding alleen verschuift
$D$ met ongeveer \qty{0.1}{km}; met echte gegevens zou de ruis van elk station groter zijn,
en zou de \omterm{def:b1:earthquake-source:locking-depth}{vastklemdiepte} tegelijk op alle stations worden gefit.

\textbf{22.} De \omterm{def:b1:earthquake-source:locking-depth}{vastklemdiepte}, \qty{15}{km}, is gelijk aan de diepte van de diepste
aardbevingen. Het vastgeklemde deel van de breuk is haar brosse deel, de \omterm{def:b1:stress-strain-rheology:strength-envelope}{seismogene laag},
waar wrijving haar tussen aardbevingen vasthoudt en in aardbevingen breekt; daaronder
kruipt de warmere rots (\cref{ch:b1:stress-strain-rheology}) en glijdt de breuk gestaag,
dat is de verschuiving met de snelheid $V$ onder $D$ van het model.

\textbf{23.} $V/(\pi D) = \num{0.030}/(\pi\times\num{15000}) = \num{6.4e-7}$ per jaar.

\textbf{24.} Met $M_0 = \qty{1.056e21}{N.m}$ is $\log_{10}M_0 = 21.024$ en
$M_w = \tfrac23(21.024 - 9.1) = 7.949$. Op 0.01 nauwkeurig heeft de aardbeving die het
hele tekort van het segment zou vrijgeven de \omterm{def:b1:earthquake-source:mw}{momentmagnitude} \textbf{7.95}.
\end{solution}
